\originalchapter{牛顿定律与常见力}\label{ch:newton}
\originalsection{惯性系与运动方程}
物体无合力时保持匀速直线运动, 这是牛顿第一定律, 同时界定了惯性参考系. 相对某惯性系作匀速平动的系仍为惯性系; 转动或加速的系一般不是. 在惯性系中, 对质量恒定的质点,
\begin{equation}
\sum\vct F=m\vct a.
\end{equation}
左边必须是作用在所选物体上的所有外力, 右边是同一参考系中的加速度. 质量恒定这一限制在第\ref{ch:massflow}章变质量问题中必须重新处理. 牛顿第二定律把相互作用与运动的变化联系起来, 不是“运动物体一定受到沿运动方向的力”.

第三定律说两物体之间的作用力等大反向: $\vct F_{A\to B}=-\vct F_{B\to A}$. 两力作用在不同物体上, 所以在单个物体的受力图中不会自行抵消. 对整个系统求总动量时, 内力的成对抵消很有用. 万有引力和理想接触满足这种形式; 电磁场带动量时需把场纳入系统, 本册不在粒子系统中滥用此抵消.

受力图先选研究对象并画真实力, 再选方便的轴. 约束决定运动形式, 如静止或沿斜面; 由投影方程求未知力. “静止”只能给 $\vct a=0$, 不能提前把未知摩擦力写成最大值.

\originalsection{引力、支持力、绳与弹簧}
两质点间引力为 $\vct F=-GmM\uv r/r^2$, $\vct r$ 从吸引中心指向被吸引物体. 球对称天体外部同样等效于质量集中于球心; 推导见第\ref{ch:gravity}章. 地表小高度范围用 $m\vct g$ 近似, $\vct g$ 向下. 支持力 $N$ 是接触面阻止互相穿透的约束力, 垂直光滑接触面. 它不必等于 $mg$, 且普通接触不能“拉住”物体, 故 $N\ge0$; 算出 $N<0$ 表示已失去接触.

绳张力沿绳方向拉物体. 理想绳无质量且不可伸长, 理想滑轮无质量无摩擦时绳两端张力相等; 有质量绳的微段需靠两端张力差加速, 张力就会变化. 对线性弹簧, 伸长 $x$ 相对自然长度, 恢复力 $F=-kx$, $k>0$ 是劲度系数. 该关系限于弹性线性范围, 不能用于任意大形变.

\originalsection{摩擦与斜面}
静摩擦在未发生相对滑动时满足 $|f_s|\le\mu_sN$, 方向由阻止相对滑动的需要确定. 动摩擦常用近似 $|f_k|=\mu_kN$, 与相对滑动速度反向, 不是总与质心速度反向. 理想无滑滚动也可能需要静摩擦来改变角速度.

斜面倾角 $\beta$ 下, 静止滑块的法向平衡给 $N=mg\cos\beta$, 切向平衡需 $f_s=mg\sin\beta$ 向上. 能静止的条件为 $\tan\beta\le\mu_s$. 若已向下滑, 加速度 $a=g(\sin\beta-\mu_k\cos\beta)$. 若滑块初始向上, 动摩擦改为向下, 加速度沿上坡正轴为 $-g(\sin\beta+\mu_k\cos\beta)$; 停住后须重新判断是否会静止或倒滑.

\begin{example}
在光滑水平面上, 质量 $m$ 的物体作半径 $R$ 的圆周运动, 绳通过圆心牵引. 求张力. 若改为圆锥摆, 绳与竖直夹角 $\beta$, 绳长 $l$, 求角速度.
\end{example}
\begin{solution}
水平圆周中张力为径向合力, $T=mv^2/R$. 圆锥摆的竖直方向平衡 $T\cos\beta=mg$, 水平向心分量 $T\sin\beta=m\omega^2l\sin\beta$. 非零摆角时约去 $\sin\beta$, 得 $\omega^2=g/(l\cos\beta)$. 两个分量共用同一个张力; 不能把 $T=mg$ 与水平向心关系同时使用.
\end{solution}
竖直圆周的绳还要求 $T\ge0$. 在最高点, 指向圆心向下, 有 $T+mg=mv^2/R$. 因而绳保持绷紧的最低顶点速率为 $\sqrt{gR}$. 此条件和能量守恒合用才给出从底部出发完成圆周的最小速率; 刚性杆可以推和拉, 因而阈值不同.
